\section{Coordinate Consequences}

The Euclidean scale-shape theorem identifies the uniform component of the
support register with regular scale and the zero-mean residual with
non-homothetic support variation. This section records the resulting
coordinates.

Let
\[
\mathbf h
=
\bar h\,\mathbf 1+\mathbf r
\in\mathcal C_D,
\qquad
\sum_{i=1}^{6}r_i=0.
\]

\subsection{Mean support and residual}

The mean support is
\[
\bar h
=
\frac{1}{6}\sum_{i=1}^{6}h_i.
\]

Relative to the reference regular dodecahedron
\[
D=P(h_0\mathbf 1),
\]
the associated regular scale ratio is
\[
\lambda(\mathbf h)
=
\frac{\bar h}{h_0}.
\]

Hence
\[
P(\bar h\,\mathbf 1)
=
\lambda(\mathbf h)D.
\]

The shape residual is
\[
\mathbf r
=
\mathbf h-\bar h\,\mathbf 1,
\]
with coordinates
\[
r_i=h_i-\bar h.
\]

Adding a common amount to every support coordinate leaves the residual
unchanged:
\[
(\mathbf h+a\mathbf 1)
-
\overline{(\mathbf h+a\mathbf 1)}\,\mathbf 1
=
\mathbf r.
\]

Under positive scaling,
\[
\mathbf h\longmapsto c\mathbf h,
\]
one has
\[
\bar h\longmapsto c\bar h
\qquad\text{and}\qquad
\mathbf r\longmapsto c\mathbf r.
\]

\subsection{Shape magnitude}

Define
\[
\sigma(\mathbf h)
=
\|\mathbf r\|
=
\left(
\sum_{i=1}^{6}(h_i-\bar h)^2
\right)^{1/2}.
\]

This is the Euclidean distance from \(\mathbf h\) to the uniform subspace:
\[
\sigma(\mathbf h)
=
\operatorname{dist}(\mathbf h,U).
\]

\begin{corollary}[Regularity criterion]
For every \(\mathbf h\in\mathcal C_D\),
\[
\sigma(\mathbf h)=0
\]
if and only if \(P(\mathbf h)\) is regular.
\end{corollary}

The quantity \(\sigma\) measures support-space departure from the regular
ray. It is not asserted to be a mechanical strain, energy, or physical
deformation cost.

\subsection{Shape direction}

When
\[
\sigma(\mathbf h)>0,
\]
define the normalized residual
\[
\widehat{\mathbf r}
=
\frac{\mathbf r}{\|\mathbf r\|}.
\]

Then
\[
\widehat{\mathbf r}\in H,
\qquad
\|\widehat{\mathbf r}\|=1.
\]

Since \(\dim H=5\), normalized shape directions lie on
\[
S(H)
=
\left\{
\mathbf v\in H:
\|\mathbf v\|=1
\right\}
\cong S^4.
\]

Every nonregular support register therefore has the unique form
\[
\mathbf h
=
\bar h\,\mathbf 1
+
\sigma\widehat{\mathbf r}.
\]

Thus a nonregular state is specified by

\begin{enumerate}
\item mean support \(\bar h\);
\item shape magnitude \(\sigma\); and
\item normalized shape direction \(\widehat{\mathbf r}\in S^4\).
\end{enumerate}

At a regular state, \(\sigma=0\), and no shape direction is defined.

\subsection{Scale-free shape}

Because
\[
\sigma(c\mathbf h)=c\sigma(\mathbf h),
\]
the magnitude \(\sigma\) depends on absolute scale.

Define the relative shape magnitude
\[
\eta(\mathbf h)
=
\frac{\sigma(\mathbf h)}{\bar h}.
\]

Since \(\bar h>0\) on the positive support domain, \(\eta\) is well defined.
It is invariant under positive scaling:
\[
\eta(c\mathbf h)
=
\eta(\mathbf h).
\]

Define also the normalized support register
\[
\mathbf q
=
\frac{\mathbf h}{\bar h}.
\]

Then
\[
\frac{1}{6}\sum_{i=1}^{6}q_i=1
\]
and
\[
\mathbf q
=
\mathbf 1+\eta\widehat{\mathbf r}.
\]

Equivalently,
\[
\mathbf h
=
\bar h
\left(
\mathbf 1+\eta\widehat{\mathbf r}
\right).
\]

The factor \(\bar h\) records absolute scale. The pair
\[
(\eta,\widehat{\mathbf r})
\]
records scale-free support shape.

\subsection{Orthogonal identity}

The canonical decomposition is orthogonal, so
\[
\|\mathbf h\|^2
=
\|\bar h\,\mathbf 1\|^2+\|\mathbf r\|^2.
\]

Since
\[
\|\mathbf 1\|^2=6,
\]
this becomes
\[
\|\mathbf h\|^2
=
6\bar h^2+\sigma^2.
\]

After division by \(\bar h^2\),
\[
\frac{\|\mathbf h\|^2}{\bar h^2}
=
6+\eta^2.
\]

These are identities in the Euclidean metric on support space. They are not
formulas for volume, surface area, energy, or any other physical quantity.

\subsection{Symmetry behavior}

For every rotational symmetry \(g\),
\[
\overline{\rho(g)\mathbf h}
=
\bar h.
\]

Because \(\rho(g)\) is orthogonal,
\[
\sigma(\rho(g)\mathbf h)
=
\sigma(\mathbf h)
\]
and
\[
\eta(\rho(g)\mathbf h)
=
\eta(\mathbf h).
\]

The normalized direction transforms equivariantly:
\[
\widehat{\mathbf r}(\rho(g)\mathbf h)
=
\rho(g)\widehat{\mathbf r}(\mathbf h).
\]

Thus \(\bar h\), \(\sigma\), and \(\eta\) are scalar symmetry invariants,
whereas \(\widehat{\mathbf r}\) retains the oriented distribution of support
variation among the six face pairs.

\subsection{Shape space and basis choice}

The five-dimensional subspace
\[
H
=
\left\{
\mathbf r\in\mathbb R^6:
\sum_i r_i=0
\right\}
\]
is canonical. An ordered basis within it is not.

For example,
\[
\begin{aligned}
\mathbf b_1&=(1,-1,0,0,0,0),\\
\mathbf b_2&=(0,1,-1,0,0,0),\\
\mathbf b_3&=(0,0,1,-1,0,0),\\
\mathbf b_4&=(0,0,0,1,-1,0),\\
\mathbf b_5&=(0,0,0,0,1,-1)
\end{aligned}
\]
span \(H\), but depend on the chosen ordering of the opposite-face pairs.

A preferred basis would require an additional criterion. No such choice is
made here.

\subsection{Summary}

For every support register in \(\mathcal C_D\),
\[
\mathbf h
=
\bar h\,\mathbf 1+\mathbf r.
\]

For nonregular states,
\[
\mathbf h
=
\bar h\,\mathbf 1
+
\sigma\widehat{\mathbf r}
=
\bar h
\left(
\mathbf 1+\eta\widehat{\mathbf r}
\right).
\]

The resulting coordinates separate

\begin{itemize}
\item absolute scale through \(\bar h\);
\item absolute shape departure through \(\sigma\);
\item scale-free shape departure through \(\eta\); and
\item shape direction through \(\widehat{\mathbf r}\in S^4\).
\end{itemize}

Section~6 applies these coordinates to representative support registers.
